\documentclass[11pt,a4paper]{article}

\usepackage[T1]{fontenc}
\usepackage[utf8]{inputenc}
\usepackage[brazil]{babel}
\usepackage[margin=2.4cm]{geometry}
\usepackage{microtype}
\usepackage{booktabs}
\usepackage{tabularx}
\usepackage[hidelinks]{hyperref}

\setlength{\parindent}{0pt}
\setlength{\parskip}{6pt}
\renewcommand{\arraystretch}{1.15}
\newcommand{\cmd}[1]{\texttt{#1}}

\begin{document}

\begin{center}
    {\small PUCRS \enspace $\cdot$ \enspace Sistemas Operacionais \enspace $\cdot$ \enspace Trabalho Prático 2}\\[6pt]
    {\LARGE\bfseries Manual do usuário}\\[4pt]
    {\large Simulador de gerenciamento de memória}\\[8pt]
    Jayme Fortes, Rodrigo Gagliardi e Lucca Marques
\end{center}

\section{Instalação e execução}

O simulador lê um arquivo de texto com pedidos de alocação, liberação e
acesso à memória. É executado no terminal e requer um \textbf{JDK 17 ou
superior}. A compilação usa \cmd{--release 17} e a biblioteca padrão do Java.

Extraia o ZIP e abra o terminal na raiz do projeto, onde ficam as pastas
\cmd{src/} e \cmd{tests/}. Para confirmar se o JDK está no \cmd{PATH},
execute \cmd{javac -version} e \cmd{java -version}.

\subsection*{Windows com PowerShell}

Execute o atalho abaixo para compilar e abrir o simulador no modo interativo.

\begin{verbatim}
.\start
\end{verbatim}

O atalho usa o JDK local em \cmd{oracleJdk-27/}, quando disponível, ou
o JDK do \cmd{PATH}. Com o JDK no \cmd{PATH}, também é possível compilar
e iniciar manualmente no PowerShell.

\begin{verbatim}
$fontes = (Get-ChildItem src -Recurse -File -Filter '*.java').FullName
javac -encoding UTF-8 --release 17 -d out $fontes
java -cp out memory.Main
\end{verbatim}

\subsection*{Linux e macOS}

Execute os comandos abaixo a partir da raiz do projeto.

\begin{verbatim}
javac -encoding UTF-8 --release 17 -d out $(find src -name '*.java')
java -cp out memory.Main
\end{verbatim}

Continue para a execução após uma compilação bem-sucedida. Se o terminal
informar que \cmd{javac} não foi encontrado, verifique se o JDK está
instalado e se a pasta \cmd{bin} dele está no \cmd{PATH}.

\clearpage

\section{Parâmetros e exemplos}

Após iniciar o programa sem argumentos, escolha o modo e o tamanho
da memória. Informe uma potência de 2, como 32 ou 64. O modo 1 pede
a política. O modo 3 pede o tamanho da página, que também deve ser
potência de 2 e menor ou igual ao tamanho da memória. Em seguida,
informe o arquivo de requisições e escolha a forma de execução.

\begin{itemize}
    \item \cmd{s} ativa o passo a passo. Pressione Enter para avançar
    após a exibição de cada comando e do estado da memória.
    \item \cmd{n} executa o arquivo inteiro e imprime os resultados sem pausas.
\end{itemize}

O caminho do arquivo pode ser relativo à raiz do projeto, como
\cmd{tests/example-1.txt}, ou absoluto. Para arquivos em \cmd{tests/},
também é possível digitar um nome curto, como \cmd{example-1}.

\subsection*{Exemplos interativos}

Digite as respostas da tabela na ordem indicada, uma por vez.
Todos os exemplos usam o modo contínuo.

\begin{center}
\small
\begin{tabularx}{\linewidth}{@{}lX@{}}
\toprule
Modo & Respostas em ordem \\
\midrule
1 --- Partições variáveis & \cmd{1}, \cmd{32}, \cmd{1} (Worst-Fit),
\cmd{tests/example-1.txt}, \cmd{n}. Para Circular-Fit, escolha
\cmd{2} na pergunta da política. \\
2 --- Sistema buddy & \cmd{2}, \cmd{64}, \cmd{tests/example-3.txt}, \cmd{n}. \\
3 --- Paginação simples & \cmd{3}, \cmd{64}, \cmd{8} (tamanho da página),
\cmd{tests/example-4.txt}, \cmd{n}. \\
\bottomrule
\end{tabularx}
\end{center}

Com o arquivo \cmd{example-1.txt}, o resumo esperado é de seis admitidos
e um rejeitado no Worst-Fit e de sete admitidos e zero rejeitados no
Circular-Fit. Os exemplos de buddy e paginação acima dão quatro admitidos
e um rejeitado cada. Na paginação, \cmd{ACC(B,8)} resulta no endereço
físico 32.

\subsection*{Execução direta por argumentos}

Após compilar, é possível fornecer todos os parâmetros na linha de
comando. Este exemplo executa partições variáveis com Worst-Fit, 64
posições de memória e o arquivo \cmd{tests/teste1.txt}.

\begin{small}
\begin{verbatim}
java -cp out memory.Main tests/teste1.txt --mode 1 --mem 64 --policy worst --continuous
\end{verbatim}
\end{small}

O resumo esperado é de sete admitidos e sete rejeitados. Em
\cmd{--policy}, use \cmd{worst} ou \cmd{circular}. Para paginação,
informe \cmd{--mode 3} e \cmd{--page} com o tamanho da página.
Use \cmd{--step} para avançar com Enter.

\clearpage

\section{Arquivo de requisições}

Salve o arquivo em UTF-8, com um comando por linha. Os endereços lógicos
de um processo começam em zero. Linhas em branco são ignoradas.

\begin{center}
\small
\begin{tabularx}{\linewidth}{@{}lX@{}}
\toprule
Comando & Efeito \\
\midrule
\cmd{IN(ID,tamanho)} & Solicita a alocação do número indicado de posições.
O tamanho deve ser inteiro e positivo. \\
\cmd{OUT(ID)} & Libera a memória ocupada pelo processo. \\
\cmd{ACC(ID,endereço)} & Traduz o endereço lógico no modo 3.
Nos modos 1 e 2, exibe um aviso e segue a execução. \\
\bottomrule
\end{tabularx}
\end{center}

Por exemplo, o arquivo abaixo aloca 6 posições para \cmd{A}, acessa
o último endereço lógico válido e depois libera o processo.

\begin{verbatim}
IN(A,6)
ACC(A,5)
OUT(A)
\end{verbatim}

\section{Como interpretar a saída}

A saída mostra o estado inicial e o resultado de cada comando. Nos modos
1 e 2, o mapa usa intervalos como \cmd{[0-7]}, com os dois extremos
incluídos. O mapa indica quais regiões estão livres e qual processo ocupa
cada região alocada. No modo 1, a lista mostra o total livre e o maior
bloco livre. No buddy, a saída também informa a fragmentação interna total.

Na paginação, o mapa lista o dono de cada frame e os frames disponíveis.
Uma entrada \cmd{p1->f4} associa a página 1 ao frame 4. O comando
\cmd{ACC} mostra a página, o deslocamento, o frame e o endereço físico.
A saída também indica a fragmentação interna total e informa quando a
tabela de páginas é liberada.

O resumo final apresenta o total de processos admitidos, os rejeitados por
falta de espaço, a ocupação máxima e a fragmentação média.

\subsection*{Mensagens comuns}

\begin{itemize}
    \item \textbf{Espaço insuficiente de memória.} O pedido não cabe no
    estado atual. O processo fica fora da memória.
    \item \textbf{Processo já existe.} O ID do \cmd{IN} já está em uso.
    \item \textbf{Processo não encontrado.} O ID informado em \cmd{OUT}
    ou em um \cmd{ACC} de paginação está ausente.
    \item \textbf{Endereço inválido.} No modo 3, o endereço é negativo
    ou alcança ou ultrapassa o tamanho do processo.
    \item \textbf{Linha inválida.} O comando foi ignorado. Verifique seu formato.
\end{itemize}

\end{document}
